\scanpage{81}
% archive foliation: 80
% lines 1--6
donnant naissance à
\[\widetilde A_\nu\longrightarrow\unsure{\mathfrak S_{\nu+1}|\mathcal S_{\nu+1}}\times\widehat Z^*\]
qui est \uline{surjectif}. Le noyau de
$\widetilde A_\nu\to\unsure{\mathfrak S_{\nu+1}|\mathcal S_{\nu+1}}$
est noté $A_\nu$, celui de
$\widetilde A_\nu\to\unsure{\mathfrak S_{\nu+1}|\mathcal S_{\nu+1}}\times\widehat Z^*$
est noté $A_\nu^0$. D'ailleurs on a
un homomorphisme injectif
\[\pi\longrightarrow A_\nu^0\subset A_\nu\]
% lines 7--10 with interlinear addition
via automorphismes intérieurs. Le conoyau
\inserted{de $\pi\longrightarrow A_\nu^0$}
\inserted{$(\pi\to A_\nu,\pi\to\widetilde A_\nu\ \unsure{:}\ \Gamma_\nu,\widetilde\Gamma_\nu)$}
est \deleted{$\unsure{\pi}$} noté $\Gamma_\nu^0$, donc on a des
suites exactes
% lines 11--14
\[\left\{\begin{array}{l}
1\mathbin{-}\pi\mathbin{-}A_\nu^0\longrightarrow\Gamma_\nu\longrightarrow1\\
\deleted{1\mathbin{-}\pi\mathbin{-}}\\
1\mathbin{-}A_\nu\mathbin{-}\widetilde A_\nu\longrightarrow\unsure{\mathfrak S_{\nu+1}|\mathcal S_{\nu+1}}\mathbin{-}1\\
1\unsure{\Longrightarrow|\longrightarrow}A_\nu^0\longrightarrow A_\nu\mathbin{-}\widehat Z^*\mathbin{-}1
\end{array}\right.\]
% line 15
\[\overbrace{\pi\quad\Gamma_\nu^0\quad\widehat Z^*\quad\unsure{\mathfrak S_{\nu+1}|\mathcal S_{\nu+1}}}^{\unsure{\widetilde A^\nu|\widetilde A_\nu}}\qquad\text{filtration}\]
% lines 16--19
\uline{Exemples} \textcircled{1} Soit $X$ droite projective sur
$k$ alg. clos de car. $0$, $I\subset X(k)$, $\operatorname{card}I=\nu+1$
\[T=\unsure{T_\bullet(k)|T_0(k)}=\lim_{\longleftarrow\,n}\mu_n(k),\qquad\pi=\pi_1(X\setminus I,x)\]
où $x\in X(k)-I$. Il \deleted{\unsure{y}} y a par Kummer
